%% toc: P1 2 2.1
\section{Membership, set operations, and bounded set formation}
\label{sec:p01-w02-p01-005-set-operations}

\Cref{sec:set-prerequisites-01} defines inclusion, equality by mutual
inclusion, the power set and classes, and warns that set-builder notation does
not by itself produce a set. This section supplies the rest of the elementary
vocabulary used in the chapter: membership as the primitive relation, bounded
set formation, binary union, binary intersection, relative complement and the
intersection of a nonempty set. Its purpose is to fix, for each operation, the
set-existence principle that licenses it and the exact membership condition it
satisfies, so that every later identity between sets can be proved by one
procedure.

The input is the ambient classical set theory that the chapter assumes in the
metalanguage \cite[ch.~``Sets, Relations, Functions, and Cardinality'',
opening paragraph]{FrozenFormalLogicSetOperations}. The set-existence
principles invoked below are exactly extensionality, the empty-set axiom,
pairing, union, power set and the separation schema, in the forms stated in
\cref{sec:zf-zfc-01} \cite[ch.~``The Axioms of ZF and ZFC'', sec.~``The
language and the axiom schemas'']{FrozenFormalLogicSetOperations}; each use is
named where it occurs. The conclusions are the existence, uniqueness and
membership condition of every operation
(\cref{thm:p01-w02-p01-005-separated-set,thm:p01-w02-p01-005-membership}),
a bounding set for Cartesian products
(\cref{thm:p01-w02-p01-005-product-bound}), the algebraic laws of the
operations (\cref{thm:p01-w02-p01-005-laws}), the Boolean algebra of subsets
of a set (\cref{thm:p01-w02-p01-005-subset-algebra}), and the failure of the
operations that have no bounding set (\cref{thm:p01-w02-p01-005-no-universal}).
The operations other than those named in the axioms, the lemmas and their
proofs are editorial completions built from the axioms just listed.

\subsection{Membership}
\label{sec:p01-w02-p01-005-membership}

\begin{hbconvention}{def:p01-w02-p01-005-membership}{Membership as the primitive relation}
The ambient set theory is a one-sorted first-order theory in the language
\(\Lang_{\in}=\{\in\}\): equality is logical identity and \(\in\) is a binary
relation symbol \cite[ch.~``The Axioms of ZF and ZFC'', paragraph before
sec.~``The language and the axiom schemas'']{FrozenFormalLogicSetOperations}.
\begin{enumerate}
  \item The variables range over sets. The formula \(x\in a\) is read
        ``\(x\) is a member of \(a\)'' and \(x\notin a\) abbreviates
        \(\neg\,x\in a\).
  \item Membership is not defined. The other notions of this section and of
        \cref{sec:set-prerequisites-01} are introduced by definition from
        \(\in\) and \(=\): either as an
        abbreviation of an \(\Lang_{\in}\)-formula, as inclusion
        \(A\subseteq B\) is in \cref{sec:set-prerequisites-01}, or as a name
        for the unique set whose existence and uniqueness have been proved, as
        each operation below is.
  \item Two sets are equal exactly when they have the same members. From
        right to left this is the axiom of extensionality; from left to right
        it is the substitutivity of identity in the underlying logic. The
        clause ``equal iff each is a subset of the other'' of
        \cref{sec:set-prerequisites-01} is this equivalence with the definition
        of \(\subseteq\) unfolded.
  \item For a class \(C=\{x:\varphi(x,\bar p)\}\) in the sense of
        \cref{sec:set-prerequisites-01}, \(x\in C\) abbreviates
        \(\varphi(x,\bar p)\). Such an abbreviation adds no class variable to
        \(\Lang_{\in}\) and asserts no existence
        \cite[ch.~``The Axioms of ZF and ZFC'', convention
        ``Classes'']{FrozenFormalLogicSetOperations}.
\end{enumerate}
\end{hbconvention}

\begin{hbconvention}{def:p01-w02-p01-005-existence-principles}{Set-existence principles used in this section}
The following sentences and schema of \(\Lang_{\in}\) are the only
set-existence and identity principles this section and its proofs use. They
are stated in \cref{sec:zf-zfc-01} \cite[ch.~``The Axioms of ZF and ZFC'',
definitions ``Extensionality'', ``Empty set'', ``Pairing, union, and power
set'' and ``Separation schema'']{FrozenFormalLogicSetOperations}.
\begin{enumerate}
  \item \emph{Extensionality:}
        \(\forall a\,\forall b\,
          \bigl(\forall x\,(x\in a\leftrightarrow x\in b)\to a=b\bigr)\).
  \item \emph{Empty set:} \(\exists a\,\forall x\,(x\notin a)\).
  \item \emph{Pairing:} \(\forall a\,\forall b\,\exists c\,\forall x\,
          \bigl(x\in c\leftrightarrow(x=a\lor x=b)\bigr)\).
  \item \emph{Union:} \(\forall a\,\exists u\,\forall x\,
          \bigl(x\in u\leftrightarrow\exists y\,(y\in a\land x\in y)\bigr)\).
  \item \emph{Power set:} \(\forall a\,\exists p\,\forall x\,
          \bigl(x\in p\leftrightarrow x\subseteq a\bigr)\), where
          \(x\subseteq a\) abbreviates \(\forall z\,(z\in x\to z\in a)\).
  \item \emph{Separation:} for every formula \(\varphi(x,\bar p)\) of
        \(\Lang_{\in}\) in which \(b\) is not free,
        \[
          \forall a\,\exists b\,\forall x\,
          \bigl(x\in b\leftrightarrow(x\in a\land\varphi(x,\bar p))\bigr),
        \]
        where the parameters \(\bar p\) are arbitrary sets.
\end{enumerate}
Each existence principle asserts that some set has a membership condition;
extensionality makes that set unique. Separation is the only principle among
them that takes an arbitrary condition, and it yields only a subset of a set
\(a\) given in advance. No principle here asserts that \(\{x:\varphi(x)\}\) is
a set for an arbitrary \(\varphi\).
\end{hbconvention}

\subsection{Bounded set formation}
\label{sec:p01-w02-p01-005-bounded}

\begin{hbdefinition}{def:p01-w02-p01-005-bounded-formation}{Bounded set formation}
Let \(A\) be a set, let \(\varphi(x,\bar p)\) be a formula of \(\Lang_{\in}\)
whose free variables are among \(x,\bar p\), and let \(\bar p\) be assigned
sets. The \emph{set of members of \(A\) satisfying \(\varphi\)} is written
\[
  \{x\in A:\varphi(x,\bar p)\}
\]
and denotes the unique set \(b\) such that
\[
  \forall x\,\bigl(x\in b\leftrightarrow(x\in A\land\varphi(x,\bar p))\bigr).
\]
The side conditions are the following.
\begin{enumerate}
  \item \emph{Bound.} \(A\) is a set already shown to exist. A class may not
        stand in this place.
  \item \emph{Definite condition.} The condition is a formula of
        \(\Lang_{\in}\), possibly after unfolding defined notation, with set
        parameters. A description that is not so expressed does not instantiate
        the schema.
  \item \emph{Fresh witness.} The variable for the formed set does not occur
        free in \(\varphi\). \Cref{thm:p01-w02-p01-005-no-universal}(5)
        shows that the existence claim fails without this condition.
  \item \emph{Bound variable.} The \(x\) in \(\{x\in A:\dots\}\) is bound;
        renaming it to a variable that does not occur in \(\varphi\) gives the
        same set (\cref{thm:p01-w02-p01-005-separated-set}).
\end{enumerate}
Licence: the separation schema \cite[ch.~``The Axioms of ZF and ZFC'',
definition ``Separation schema'']{FrozenFormalLogicSetOperations}, whose
instance for \(\varphi\) asserts the existence, and extensionality, which
gives the uniqueness (\cref{thm:p01-w02-p01-005-separated-set}).
\end{hbdefinition}

\begin{hblemma}{thm:p01-w02-p01-005-separated-set}{proved=annex:p01-w02-p01-005-separated-set}{Existence, uniqueness and inclusion of a separated set}
Under the side conditions of \cref{def:p01-w02-p01-005-bounded-formation},
there is exactly one set \(b\) with
\(\forall x\,(x\in b\leftrightarrow(x\in A\land\varphi(x,\bar p)))\), and this
set satisfies \(b\subseteq A\). If \(y\) is a variable that does not occur in
\(\varphi(x,\bar p)\) and \(\varphi(y,\bar p)\) is the result of replacing the
free occurrences of \(x\) by \(y\), then
\(\{y\in A:\varphi(y,\bar p)\}=\{x\in A:\varphi(x,\bar p)\}\).
\end{hblemma}

\begin{hbmethod}{def:p01-w02-p01-005-bounding-procedure}{Showing that a described collection is a set}
Input: a description of an intended collection by a condition \(\varphi\).
Output: either a proof that a set has exactly the intended members, or the
statement of what is missing.
\begin{enumerate}
  \item Find a set \(B\), obtained from sets already available by the
        empty-set, pairing, union or power-set axiom or by an earlier
        construction, together with a proof that every intended member lies in
        \(B\).
  \item Write the condition as an \(\Lang_{\in}\)-formula \(\varphi(x,\bar p)\)
        with set parameters, unfolding all defined notation.
  \item Form \(\{x\in B:\varphi(x,\bar p)\}\) by
        \cref{thm:p01-w02-p01-005-separated-set}.
  \item Prove by extensionality that this set is the intended collection:
        its members are the members of \(B\) satisfying \(\varphi\), and by the
        first step these are all the objects satisfying \(\varphi\).
\end{enumerate}
If no bounding set is found, the procedure yields nothing: the collection is
then only a class, and \cref{thm:p01-w02-p01-005-no-universal} shows that for
some conditions no bounding set exists. Inclusion \(B\subseteq C\) of two sets
is proved by showing \(x\in B\to x\in C\) for arbitrary \(x\); identity
\(B=C\) is proved by showing \(x\in B\leftrightarrow x\in C\) for arbitrary
\(x\) and applying extensionality.
\end{hbmethod}

\subsection{The operations}
\label{sec:p01-w02-p01-005-operations-defs}

The empty set \(\varnothing\), the unordered pair \(\{a,b\}\), the singleton
\(\{a\}=\{a,a\}\), the union \(\bigcup a\) of a set \(a\) and the power set
\(\mathcal P(a)\) are the unique witnesses of the empty-set, pairing, union and
power-set axioms \cite[ch.~``The Axioms of ZF and ZFC'', definitions ``Empty
set'' and ``Pairing, union, and power set'']{FrozenFormalLogicSetOperations}.
Their membership conditions are
\[
  x\notin\varnothing,\qquad
  x\in\{a,b\}\leftrightarrow(x=a\lor x=b),\qquad
  x\in\textstyle\bigcup a\leftrightarrow\exists y\,(y\in a\land x\in y),
\]
and \(x\in\mathcal P(a)\leftrightarrow x\subseteq a\). The power set of
\cref{sec:set-prerequisites-01}, written there with unbounded set-builder
notation, is a set because of the power-set axiom; it is not obtained by
separation.

\begin{hbdefinition}{def:p01-w02-p01-005-operations}{Union, intersection, relative complement and finite listing}
For sets \(A\), \(B\), \(a\), \(b\), \(c\) and \(s\):
\begin{enumerate}
  \item the \emph{union} of \(A\) and \(B\) is
        \(A\cup B\defeq\bigcup\{A,B\}\);
  \item the \emph{intersection} of \(A\) and \(B\) is
        \(A\cap B\defeq\{x\in A:x\in B\}\);
  \item the \emph{relative complement} of \(B\) in \(A\), also called the
        \emph{difference} of \(A\) and \(B\), is
        \(A\setminus B\defeq\{x\in A:x\notin B\}\);
  \item the \emph{three-element listing} is
        \(\{a,b,c\}\defeq\{a,b\}\cup\{c\}\), and the \emph{successor} of
        \(s\) is \(s\cup\{s\}\);
  \item \(A\) and \(B\) are \emph{disjoint} iff \(A\cap B=\varnothing\).
\end{enumerate}
In clauses 2 and 3 the condition has the parameter \(B\) and the bound is
\(A\), so the side conditions of \cref{def:p01-w02-p01-005-bounded-formation}
hold. In clause 1 no bound is available in advance; the union axiom supplies
the set.
\end{hbdefinition}

\begin{hbdefinition}{def:p01-w02-p01-005-big-intersection}{Intersection of a nonempty set}
For a set \(F\) with \(F\neq\varnothing\), the \emph{intersection} of \(F\) is
\[
  \textstyle\bigcap F\defeq
  \{x\in\bigcup F:\forall Y\,(Y\in F\to x\in Y)\}.
\]
The bound is \(\bigcup F\) and the parameter is \(F\). The hypothesis
\(F\neq\varnothing\) is part of the definition: for \(F=\varnothing\) the
displayed set is \(\varnothing\), which does not have the membership condition
of \cref{thm:p01-w02-p01-005-membership}(4), and no set has that condition
(\cref{thm:p01-w02-p01-005-no-universal}(4)).
\end{hbdefinition}

\begin{hblemma}{thm:p01-w02-p01-005-membership}{proved=annex:p01-w02-p01-005-membership}{Membership conditions of the operations}
For all sets \(A,B,a,b,c,s,z\) and every nonempty set \(F\):
\begin{enumerate}
  \item \(z\in A\cup B\leftrightarrow(z\in A\lor z\in B)\);
  \item \(z\in A\cap B\leftrightarrow(z\in A\land z\in B)\);
  \item \(z\in A\setminus B\leftrightarrow(z\in A\land z\notin B)\);
  \item \(z\in\bigcap F\leftrightarrow\forall Y\,(Y\in F\to z\in Y)\);
  \item \(z\in\{a,b,c\}\leftrightarrow(z=a\lor z=b\lor z=c)\) and
        \(z\in s\cup\{s\}\leftrightarrow(z\in s\lor z=s)\);
  \item \(A\subseteq A\cup B\), \(B\subseteq A\cup B\),
        \(A\cap B\subseteq A\), \(A\setminus B\subseteq A\), and
        \(\bigcap F\subseteq Y\) for every \(Y\in F\);
  \item \(\bigcap\{A,B\}=A\cap B\).
\end{enumerate}
\end{hblemma}

\begin{hblemma}{thm:p01-w02-p01-005-product-bound}{proved=annex:p01-w02-p01-005-product-bound}{A bounding set for the Cartesian product}
For all sets \(A\) and \(B\), every ordered pair
\(\langle a,b\rangle=\{\{a\},\{a,b\}\}\) with \(a\in A\) and \(b\in B\) is a
member of \(\mathcal P(\mathcal P(A\cup B))\). Consequently
\[
  \{z\in\mathcal P(\mathcal P(A\cup B)):
    \exists a\,\exists b\,(a\in A\land b\in B\land z=\langle a,b\rangle)\}
\]
is a set whose members are exactly the ordered pairs \(\langle a,b\rangle\)
with \(a\in A\) and \(b\in B\); it is the Cartesian product \(A\times B\) of
\cref{sec:set-prerequisites-02}.
\end{hblemma}

The bounding set and the axioms used follow the architecture of the solution to
the product-construction exercise of the axiomatic set-theory chapter
\cite[ch.~``Exercise Feedback'', sec.~``Set theory and forcing'', solution to
the first exercise of ``The Axioms of ZF and ZFC'']{FrozenFormalLogicSetOperations};
the annex writes every step.

\subsection{Laws of the operations}
\label{sec:p01-w02-p01-005-laws}

\begin{hbtheorem}{thm:p01-w02-p01-005-laws}{proved=annex:p01-w02-p01-005-laws}{Laws of union, intersection and relative complement}
For all sets \(A\), \(B\), \(C\):
\begin{enumerate}
  \item \emph{Idempotence:} \(A\cup A=A\) and \(A\cap A=A\).
  \item \emph{Commutativity:} \(A\cup B=B\cup A\) and \(A\cap B=B\cap A\).
  \item \emph{Associativity:} \((A\cup B)\cup C=A\cup(B\cup C)\) and
        \((A\cap B)\cap C=A\cap(B\cap C)\).
  \item \emph{Distributivity:} \(A\cap(B\cup C)=(A\cap B)\cup(A\cap C)\) and
        \(A\cup(B\cap C)=(A\cup B)\cap(A\cup C)\).
  \item \emph{Absorption:} \(A\cup(A\cap B)=A\) and \(A\cap(A\cup B)=A\).
  \item \emph{Empty set:} \(A\cup\varnothing=A\), \(A\cap\varnothing=\varnothing\),
        \(A\setminus\varnothing=A\) and \(A\setminus A=\varnothing\).
  \item \emph{Relative De Morgan laws:}
        \(C\setminus(A\cup B)=(C\setminus A)\cap(C\setminus B)\) and
        \(C\setminus(A\cap B)=(C\setminus A)\cup(C\setminus B)\).
  \item \emph{Difference through intersection:}
        \(A\setminus B=A\setminus(A\cap B)\).
  \item \emph{Inclusion:} the four statements \(A\subseteq B\),
        \(A\cup B=B\), \(A\cap B=A\) and \(A\setminus B=\varnothing\) are
        equivalent.
  \item \emph{Complement inside a fixed set:} if \(A\subseteq C\), then
        \(C\setminus(C\setminus A)=A\), \(A\cup(C\setminus A)=C\),
        \(A\cap(C\setminus A)=\varnothing\), \(A\cup C=C\) and \(A\cap C=A\).
\end{enumerate}
\end{hbtheorem}

\begin{hbcorollary}{thm:p01-w02-p01-005-subset-algebra}{proved=annex:p01-w02-p01-005-subset-algebra}{The subsets of a set form a Boolean algebra}
For every set \(S\), the set \(\mathcal P(S)\) is closed under \(\cup\),
\(\cap\) and \(X\mapsto S\setminus X\), and contains \(\varnothing\) and \(S\).
Union and intersection are functions from \(\mathcal P(S)\times\mathcal P(S)\)
to \(\mathcal P(S)\), and complement is a function from \(\mathcal P(S)\) to
\(\mathcal P(S)\), in the sense of \cref{sec:set-prerequisites-03}. With
meet \(\cap\), join \(\cup\), complement \(X\mapsto S\setminus X\), bottom
\(\varnothing\) and top \(S\), it satisfies the bounded distributive lattice
and complement laws, so it is a Boolean algebra in the sense of
\cref{def:p01-bpi-boolean-algebra}. Its order \(X\leq Y\), defined by
\(X\cap Y=X\), is inclusion \(X\subseteq Y\), and the algebra is nontrivial iff
\(S\neq\varnothing\).
\end{hbcorollary}

\subsection{Worked example}
\label{sec:p01-w02-p01-005-example}

\begin{hbexample}{def:p01-w02-p01-005-example}{Operations on three pure sets}
Put \(e_0=\varnothing\), \(e_1=\{\varnothing\}\) and \(e_2=\{\{\varnothing\}\}\);
the empty-set and pairing axioms give all three. They are pairwise distinct.
First, \(\varnothing\in e_1\) and \(\varnothing\notin e_0\), so \(e_0\neq e_1\).
Second, \(e_1\in e_2\) and \(e_1\notin e_0\), so \(e_0\neq e_2\). Third, the
only member of \(e_2\) is \(\{\varnothing\}\), and
\(\{\varnothing\}\neq\varnothing\) because \(\varnothing\) is a member of the
first and not of the second; so \(\varnothing\notin e_2\) while
\(\varnothing\in e_1\), and \(e_1\neq e_2\). Each inequality uses the
left-to-right direction of clause 3 of
\cref{def:p01-w02-p01-005-membership}.

Let \(A=\{e_0,e_1\}\) and \(B=\{e_1,e_2\}\). By
\cref{thm:p01-w02-p01-005-membership}, for every \(z\):
\begin{itemize}
  \item \(z\in A\cup B\) iff \(z=e_0\lor z=e_1\lor z=e_2\), so
        \(A\cup B=\{e_0,e_1,e_2\}\) by extensionality.
  \item \(z\in A\cap B\) iff \((z=e_0\lor z=e_1)\land(z=e_1\lor z=e_2)\). The
        disjunct \(z=e_0\land z=e_2\) and the disjuncts that pair \(z=e_0\)
        or \(z=e_2\) with \(z=e_1\) are impossible by distinctness, so the
        condition reduces to \(z=e_1\), and \(A\cap B=\{e_1\}\).
  \item \(z\in A\setminus B\) iff \((z=e_0\lor z=e_1)\land z\neq e_1\land
        z\neq e_2\), which reduces to \(z=e_0\); so \(A\setminus B=\{e_0\}\),
        and likewise \(B\setminus A=\{e_2\}\).
  \item With \(C=A\cup B\), the bounded formation
        \(\{x\in C:x\notin A\cap B\}\) is \(C\setminus(A\cap B)=\{e_0,e_2\}\).
        The second relative De Morgan law of \cref{thm:p01-w02-p01-005-laws}
        computes the same set as
        \((C\setminus A)\cup(C\setminus B)=\{e_2\}\cup\{e_0\}\).
  \item \(A\cap\{e_2\}=\varnothing\), since \(e_2\notin A\); so \(A\) and
        \(\{e_2\}\) are disjoint, while \(A\) and \(B\) are not, since
        \(e_1\in A\cap B\).
\end{itemize}
The same computations show that \(e_1=e_0\cup\{e_0\}\) and
\(A=e_1\cup\{e_1\}\): the successor operation of
\cref{def:p01-w02-p01-005-operations} carries \(e_0\) to \(e_1\) and \(e_1\)
to \(A\). The set \(e_2\) is not the successor of \(e_1\), because
\(e_0\in A\) and \(e_0\notin e_2\).
\end{hbexample}

\subsection{Failure without a bounding set}
\label{sec:p01-w02-p01-005-failure}

\begin{hbtheorem}{thm:p01-w02-p01-005-no-universal}{proved=annex:p01-w02-p01-005-no-universal}{Collections with no bounding set}
\begin{enumerate}
  \item For every set \(a\), the set \(r_a=\{x\in a:x\notin x\}\) is not a
        member of \(a\).
  \item There is no set \(V\) such that \(\forall x\,(x\in V)\).
  \item For every set \(A\), there is no set \(D\) such that
        \(\forall x\,(x\in D\leftrightarrow x\notin A)\); the absolute
        complement \(\{x:x\notin A\}\) is a proper class.
  \item There is no set \(I\) such that
        \(\forall x\,\bigl(x\in I\leftrightarrow\forall Y\,(Y\in\varnothing\to x\in Y)\bigr)\);
        the intersection of the empty set, taken with the membership condition
        of \cref{thm:p01-w02-p01-005-membership}(4), is a proper class. The
        set \(\{x\in\bigcup\varnothing:\forall Y\,(Y\in\varnothing\to x\in Y)\}\)
        is \(\varnothing\).
  \item The sentence
        \(\forall a\,\exists b\,\forall x\,\bigl(x\in b\leftrightarrow(x\in a\land x\notin b)\bigr)\),
        which would be the separation instance for the formula \(x\notin b\)
        if the formed variable were allowed to occur free, is false: for
        \(a=\{\varnothing\}\) there is no such \(b\).
\end{enumerate}
\end{hbtheorem}

\begin{hbwarning}{def:p01-w02-p01-005-unbounded-warning}{Operations that need a bound}
Union and power set come from their own axioms; intersection, relative
complement and every separated set need a bound that is already a set. A
complement is therefore always relative to a named set \(C\), written
\(C\setminus A\), and \(\bigcap F\) is used only after \(F\neq\varnothing\)
has been shown. The class \(\{x:x\notin x\}\) of the warning in
\cref{sec:set-prerequisites-01} is the case \(\varphi(x)=x\notin x\): by
\cref{thm:p01-w02-p01-005-no-universal}(1) its bounded version
\(\{x\in a:x\notin x\}\) exists for every set \(a\) and lies outside \(a\), so
the unbounded class is not a set \cite[ch.~``The Axioms of ZF and ZFC'',
warning ``Unrestricted comprehension'']{FrozenFormalLogicSetOperations}. The
laws of \cref{thm:p01-w02-p01-005-laws} that mention a complement hold for
relative complements; there is no absolute form of them in this set theory.
\end{hbwarning}

\subsection{Relationships}
\label{sec:p01-w02-p01-005-relationships}

The section extends \cref{sec:set-prerequisites-01}: inclusion and the
mutual-inclusion form of equality are the instances of clause 3 of
\cref{def:p01-w02-p01-005-membership} and of the inclusion procedure of
\cref{def:p01-w02-p01-005-bounding-procedure}. The axioms it uses are those
of \cref{sec:zf-zfc-01}. That section writes the successor
\(S(x)=x\cup\{x\}\) in the axiom of infinity and forms \(\omega\) by separation
from an inductive set
\cite[ch.~``The Axioms of ZF and ZFC'', definitions ``Infinity'' and ``The set
\(\omega\)'']{FrozenFormalLogicSetOperations}; the successor is clause 4 of
\cref{def:p01-w02-p01-005-operations}, and the formation of \(\omega\) is an
instance of \cref{thm:p01-w02-p01-005-separated-set}.
\Cref{thm:p01-w02-p01-005-product-bound} supplies the bounding set that the
unbounded definition of \(A\times B\) in \cref{sec:set-prerequisites-02}
requires. The Boolean algebra of sets \(\mathcal P(S)\), used with filters and
ultrafilters in \cref{def:p01-bpi-filter-ideal}, is
\cref{thm:p01-w02-p01-005-subset-algebra}. The closure constructions of
\cref{sec:p01-equivalence-closure-and-quotients} form relations by separation
from \(X\times X\) and intersect nonempty families, which are
\cref{thm:p01-w02-p01-005-separated-set,thm:p01-w02-p01-005-product-bound}
and \cref{def:p01-w02-p01-005-big-intersection}.

\subsection{Sources}
\label{sec:p01-w02-p01-005-sources}

The ambient-theory statement, the subset, power-set and class
definitions, the language \(\Lang_{\in}\), the axioms of extensionality, empty
set, pairing, union, power set and separation, the class convention, the
successor notation, the definition of \(\omega\), the unrestricted
comprehension warning and the product-exercise architecture are taken from the
frozen manuscript \cite{FrozenFormalLogicSetOperations} at the locators given
above. The definitions of binary union, intersection, relative complement,
intersection of a nonempty set and three-element listing, the procedure, the
lemmas, the laws, the Boolean-algebra corollary, the failure theorem and all
their proofs are editorial completions from those axioms and are not
attributed to an external presentation.
